\section{Introduction}

In a daily card-transaction pipeline of about 50 million rows at a large US card issuer, an upstream merchant feed began sending duplicate keys. A broadcast join multiplied the affected transactions, and the post-load monitors for row counts and key uniqueness paged the on-call engineer. The fix was the obvious one: drop the duplicates. Every monitor turned green. But the duplicates were not identical. Some carried the amount before settlement adjustments and some the amount after, and an unordered de-duplication kept whichever came first. About 1,200 accounts were left with pre-adjustment amounts that no monitor could see, and the fix cost a full reprocessing of the day's feed and a one-day delay in reporting.

That incident is the problem this paper addresses. The tools that watch data quality mostly detect and alert, and the next step, letting software fix the data, is where the industry is heading: commercial platforms now advertise agents that remediate data incidents, and research prototypes let language models write cleaning code. A fix only has to make the alert go away to look successful, and the alert is the wrong judge. The same failure appears in reverse when the data is right: trimming a real promotion day back to a normal Friday also turns every monitor green, and quietly deletes real transactions. In our benchmark an agent whose fixes only had to quiet the existing monitors left 24\% of incidents with wrong data, and a rules engine that applied the obvious fix to every alert did so in 31\%. Both repeated the opening incident, keeping the pre-settlement copy of re-sent rows, in essentially every case where it occurred (Section~\ref{sec:exp}).

Automatic repair has a long and honest record of this problem. A recent study of twelve repair methods across many datasets found that repairs frequently introduce errors of their own, and asked whether the field is ready to deploy~\cite{ni2024ready}. Probabilistic systems such as HoloClean~\cite{rekatsinas2017holoclean} and learned ones such as Baran~\cite{mahdavi2020baran} report how accurate their repairs are on average, but a team deciding whether to switch one on wants something else: a promise about how often an incident will end with a wrong automatic change, and a system that hands everything else to a person.

This paper builds that promise. DataRanger is a repair agent whose central component is a \emph{gate}. Every candidate fix, whether a template or a language model's proposal, is applied to a copy and checked against evidence the fix cannot influence: the control totals that source systems publish, and the table's declared invariants. The gate applies a fix only if it matches that evidence within a threshold, and it chooses the threshold with Learn-then-Test~\cite{angelopoulos2021ltt} on past incidents with known outcomes. The result is a bound on \emph{silent harm}, the share of incidents that end with a wrong fix applied: with probability at least $1-\delta$, it stays below a target $\alpha$ that the team chooses. A stricter target never makes the system less safe; it only sends more incidents to people.

Conformal ideas have reached data cleaning before~\cite{jaeger2024cdc,biessmann2026cured,liu2026progressive}, and they guarantee something different: that the correct value lies inside a predicted set. That bounds how often a clean cell is flagged, under the assumption that test data looks like clean training data, which corrupted data violates. Our bound is on the outcome that matters in production, a wrong change reaching the table, and our evidence comes from the source rather than from the model's own confidence. We also target the unit of work data teams actually handle, an incident on a partition or a table, rather than a cell.

The guarantee rests on one assumption, that future incidents resemble the calibration incidents. We test it directly. When a new incident type appears that the calibration never saw, the bound breaks for that type, and it breaks badly (Section~\ref{sec:exp-drift}). DataRanger's repair memory addresses this with a novelty guard, a blocklist of signatures whose fixes went wrong, and recalibration in epochs.

\paragraph{What is new.} While statistical risk control via Learn-then-Test~\cite{angelopoulos2021ltt} is established in machine learning, existing data repair systems have been unable to certify a bound on wrong changes because they evaluate repairs using internal model confidence or alert suppression. Our core contribution is the formulation of an externally verifiable risk loss and the end-to-end architecture that makes data repair mathematically certifiable---a capability no prior repair system provides. The gate thresholds evidence the fix cannot influence (gaps to the source's control totals and checks of declared invariants, computed on a copy of the partition), so the same score means the same thing whichever proposer produced the fix; a score from the proposer itself does not, as the self-checking agents of Section~\ref{sec:exp-llm} show. The unit is the incident a data team handles, a partition or table, with escalation as a first-class outcome and fixes written as new versions, so an automatic decision is cheap to undo and a person's resolution becomes a label. And it runs over time: repair memory and epoch recalibration, with the conditions under which the guarantee survives them, and an empirical account of what they do when it does not. The value of calibration itself is isolated by deterministic reconciliation with the same candidates and evidence: requiring an exact match to the totals was safe but automated 32\% of incidents, against 59\% for the gate at a 1\% target; a well-chosen hand-set tolerance matched the gate on that benchmark but broke the target in every split when the source published coarser totals.

\paragraph{Contributions.}
\begin{itemize}
\item \textbf{Risk-controlled repair.} A gate that bounds the rate of wrong fixes a repair system applies, calibrated with Learn-then-Test and evidence from the source of record and declared invariants, with a guarantee for independent incidents, a cluster-level guarantee for dependent ones, and conditions under which grouping and memory preserve it (Section~\ref{sec:system}).
\item \textbf{Repair memory under drift.} Signatures, a novelty guard and epoch recalibration; a frozen-memory variant keeps the guarantee within a stable period, and the deployed variant reduced observed harm and sped recovery when the incident mix changed, without a formal guarantee (Section~\ref{sec:system}).
\item \textbf{DQ-Repair.} An open benchmark of pipeline incidents with decoys, stale control totals, a retention window and ground-truth fixes, plus a protocol on the Hospital, Flights and Beers datasets (Section~\ref{sec:bench}).
\item \textbf{Evaluation.} Over 6{,}600 incidents the gate kept silent harm at or below target in at least 95\% of calibration splits, against 24 to 31\% harm for monitor-checked and rule-based repair; the same held on incidents injected into TPC-DS; hand-tuned thresholds broke when evidence got coarser, and a self-checking agent with a hand-set tolerance exceeded the target the gate kept; open-weights language models proposed fixes that were wrong often enough to need the gate; and on the real datasets the same gate certified a bound that Baran and conformal cleaning do not offer (Section~\ref{sec:exp}).
\end{itemize}

\section{Related work}
\label{sec:related}

\paragraph{Error detection and repair.} Constraint-based repair finds a minimal change that satisfies functional dependencies or denial constraints~\cite{bohannon2005cost,chu2013holistic,cong2007improving}. HoloClean combines constraints, statistics and external signals in a probabilistic model~\cite{rekatsinas2017holoclean}. Raha and Baran learn detection and correction from a small labeling budget~\cite{mahdavi2019raha,mahdavi2020baran}, and UniClean orders heterogeneous cleaners to limit error propagation on mixed errors, with a result on the expected number of remaining errors~\cite{ding2025uniclean}. A large empirical study found that automatic repair often lowers data quality and that no method dominates~\cite{ni2024ready}. These systems repair cell values and report accuracy; none certifies a bound on wrong changes, and none uses an external source of record.

\paragraph{Language models for cleaning.} Foundation models can perform cleaning and wrangling tasks from prompts~\cite{narayan2022wrangle}, and several systems generate cleaning workflows or code with agents~\cite{autodcworkflow2025,cleaning_agents2025}. A 2026 study of agentic cleaning with risk-tiered abstention reported near-zero repair accuracy and no formal guarantee~\cite{fadlallah2026agentic}, and agents that repair data behind visualizations struggle with temporal, geographic and semantic issues~\cite{kashireddy2026viz}. We treat the model as a proposer whose output must pass a gate it cannot influence.

\paragraph{Agents for pipeline incidents.} A recent reference architecture for self-healing data pipelines combines a deterministic policy engine, LLM agents, risk-tiered approval and an incident memory, and verifies a remediation by re-running the checks that detected the failure~\cite{eshun2026selfhealing}; it reports no evaluation. Our monitor-checked baseline is exactly that verification rule, and it is the one that trims real promotion days. In earlier work we built PipeMend, a benchmark of broken pipelines in which language-model agents repair pipeline \emph{code}, and found that wrong fixes routinely passed the team's monitors~\cite{sharma2026fixorpage}. DataRanger addresses the same failure for repairs to the \emph{data} itself, and replaces an empirical finding with a certified bound.

\paragraph{Conformal and risk-controlled methods.} Conformal data cleaning flags a cell when its value falls outside a conformal prediction set and replaces it with the point prediction~\cite{jaeger2024cdc}; CURED packages it as a demonstrator~\cite{biessmann2026cured}, and a 2026 pipeline uses conformal reliability control to reject low-confidence corrections~\cite{liu2026progressive}. Learn-then-Test~\cite{angelopoulos2021ltt}, risk-controlling prediction sets~\cite{bates2021rcps} and conformal risk control~\cite{angelopoulos2024crc} certify that a tuned procedure keeps an arbitrary risk below a target. We apply Learn-then-Test to the decision a repair system makes, and test where its exchangeability assumption fails.

\paragraph{Data validation and observability.} Production systems validate data against schemas and statistics before it is used~\cite{breck2019tfdv,schelter2018deequ}, and pipeline verification compares outputs to their sources. DataRanger reuses these checks twice: as detectors, and as the evidence that decides whether a fix is applied.
